% =============================================================================
% Chapter 8 -- Testing and Early Validation         (owner: Mishal Fatima)
% All results below were produced by running the commands shown, on the
% reporting commit. Where a measurement has not been produced yet, the section
% carries a PENDING marker instead of an estimate.
% Reference: docs/dataset/08_integrity_report.md (138 checks, 0 failures).
% =============================================================================
\chapter{Testing and Early Validation}
\label{ch:testing}

Testing at this stage has two jobs: prove that what claims to be implemented
actually builds and behaves, and prove that the evaluation data which later
stages will be measured against is itself sound. The second job is the one that
is furthest along, and the first is limited by how little of the system exists.

% ---------------------------------------------------------------------------
\section{Testing Strategy}
\label{sec:test-strategy}

The strategy follows the five-layer validation plan recorded in Chapter~3, but
only the layers that can honestly be exercised now are reported:

\begin{enumerate}
  \item \textbf{Static correctness} of the front end: the compiler and the linter
        are the cheapest tests that exist, they run on every change, and they
        enforce the domain vocabulary (a fifth verdict value, or a status with no
        label mapping, is a compile error rather than a run-time surprise).
  \item \textbf{Build reproducibility}: a production build must complete and
        pre-render every route, because a demo depends on it and because a
        broken route fails here rather than on stage.
  \item \textbf{Component tests} over the surviving proof-of-concept code,
        run with \texttt{pytest}.
  \item \textbf{Data validation}: automated integrity checks over every dataset
        case (schema, commit resolution, patch application, span bounds, test
        logs, leak and secret scans).
  \item \textbf{Annotation quality}: independent double labelling with raw
        agreement and Cohen's kappa reported against a pre-declared
        threshold~\cite{cohen1960}.
\end{enumerate}

Two layers of the plan are declared but not yet runnable, and are listed as
such rather than implied: system-level acceptance tests per functional
requirement (there is no system to run them against), and comparative
benchelines against baselines (there is no model and no harness).

% ---------------------------------------------------------------------------
\section{Test Levels and Techniques}
\label{sec:test-levels}

\begin{longtable}{p{0.18\textwidth} p{0.26\textwidth} p{0.28\textwidth} p{0.22\textwidth}}
\caption{Test levels, the technique used at each level, and the current execution status.}\label{tab:test-levels}\\
\hline\textbf{Level} & \textbf{Technique} & \textbf{Tool / command} &
  \textbf{Status} \\
\hline\endfirsthead
\hline\textbf{Level} & \textbf{Technique} & \textbf{Tool / command} &
  \textbf{Status} \\
\hline\endhead
Type level & Static type checking over every route, component and mock module &
  \texttt{npm run typecheck} (\texttt{tsc --noEmit}) & Running \\
Lint & Rule-based inspection (Next.js + TypeScript rules) &
  \texttt{npm run lint} (\texttt{eslint .}) & Running \\
Build & Production compilation and static generation of all routes &
  \texttt{npm run build} (\texttt{next build}) & Running \\
Unit & Assertion-based component tests for the proof-of-concept adapters &
  \texttt{python3 -m pytest} & Present but failing to collect (D-01) \\
Data integrity & Scripted per-case checks over the dataset package &
  \texttt{docs/dataset/tools/check\_integrity.py} & Running (report committed) \\
Annotation agreement & Blind double labelling, raw agreement, Cohen's kappa,
  span overlap & \texttt{docs/dataset/tools/agreement.py} & Half complete
  (annotator A pending) \\
System / end-to-end & Scenario matrix T1--T11 from Chapter~4 &
  -- & Not started (no back end) \\
Performance & Latency and memory measurement & -- & Not started (no model) \\
\hline\end{longtable}

There is no CI pipeline in the repository yet: none of the commands above runs
automatically on a pull request. That is defect D-02 below.

% ---------------------------------------------------------------------------
\section{Test Cases}
\label{sec:test-cases}

\begin{longtable}{p{0.08\textwidth} p{0.24\textwidth} p{0.24\textwidth} p{0.20\textwidth} p{0.16\textwidth}}
\caption{Selected test cases with their expected results and observed outcomes.}\label{tab:test-cases}\\
\hline\textbf{ID} & \textbf{Method and data} & \textbf{Expected result} &
  \textbf{Actual result} & \textbf{Evidence} \\
\hline\endfirsthead
\hline\textbf{ID} & \textbf{Method and data} & \textbf{Expected result} &
  \textbf{Actual result} & \textbf{Evidence} \\
\hline\endhead
TC-01 & \texttt{npm run typecheck} over \texttt{apps/web/src} & No type errors,
  exit code 0 & Exit code 0, no diagnostics & Command output, Ch.~6 \\
TC-02 & \texttt{npm run lint} over \texttt{apps/web} & No findings, exit code 0 &
  Exit code 0, no findings & Command output, Ch.~6 \\
TC-03 & \texttt{npm run build}; full production build & Build succeeds; every
  route emitted & Succeeded; TypeScript pass in 3.2\,s; 29 of 29 static pages
  generated in 1075\,ms & Build log, Ch.~6 \\
TC-04 & Static generation of the dynamic route
  \texttt{conflicts/[id]} & One page per mock finding; no missing path &
  Eight \texttt{CONF-*} paths emitted (three named, five abbreviated in the
  build summary) & Build log \\
TC-05 & \texttt{check\_integrity.py} over all 15 pilot cases: required fields,
  base/head SHA resolution, patch SHA-256 and clean apply, span existence and
  in-bounds line ranges, content hash stability, referenced test logs, label
  leak scan, secret scan & Zero failures & 138 checks, 0 failures across 15
  cases (8--10 checks per case) &
  \texttt{docs/dataset/08\_integrity\_report.md} \\
TC-06 & \texttt{agreement.py} over both annotation sheets & Raw agreement,
  kappa per field and span overlap, or a clear statement that the data is
  incomplete & Exited with ``No case labelled by both annotators yet.''
  (annotator A has 0 of 15 verdicts filled) & Command output, below \\
TC-07 & \texttt{python3 -m pytest tests/} in \texttt{unify\_poc} & Adapter
  mapping tests pass & Collection error: \texttt{ModuleNotFoundError: No module
  named 'unify'}; the module's own imports (\texttt{adapter},
  \texttt{normalizer}) are absent from the tree & Command output; defect D-01 \\
TC-08 & Type-level check: assign a non-vocabulary string to a
  \texttt{Verdict} & Compile error & Compile error by construction
  (\texttt{VERDICTS} is a const tuple) & \texttt{src/lib/types.ts} \\
\hline\end{longtable}

% ---------------------------------------------------------------------------
\section{Test Results}
\label{sec:test-results}

\paragraph{Front end.} Three of three static and build checks pass. The build is
the strongest result available today because it exercises every route, every
import and every type in one run, and it is the check a reviewer can repeat with
one command. One advisory warning is emitted by the toolchain rather than by our
code: Node reports a deprecation warning for \texttt{url.parse()} coming from
the framework's dependencies. It does not fail the build and is recorded as
D-05 so that it is not mistaken for a clean silence.

\paragraph{Dataset.} The integrity report committed at
\texttt{docs/dataset/08\_integrity\_report.md} records 15 cases checked, 138
checks run, 0 failures, with per-case results of PASS and the check list stated
in the file. Two classes of check are explicitly \emph{not} run yet, and the
report says so: cross-split duplicate detection (no split is frozen) and a
checksum manifest of frozen cases (cases are not frozen). Reporting what the
test does not cover is part of the result.

\paragraph{Proof of concept.} The single remaining Python test module does not
collect. This is a genuine failure, not a formatting issue: the file imports
from a package that has not been laid out in the tree, and the modules it
depends on were removed from the repository. It is discussed as the edge case
below and tracked as D-01.

\paragraph{Agreement.} \textbf{[PENDING: agreement results --- run
tools/agreement.py]}

% ---------------------------------------------------------------------------
\section{Dataset Pilot Agreement}
\label{sec:agreement}

\paragraph{Protocol.} The pilot is 15 cases drawn from 3 TypeScript/React
repositories, covering all four verdicts with same-module hard negatives. Two
members label every case independently in their own sheet and may not discuss
cases until both sheets are submitted; the third member adjudicates each
disagreement into \texttt{06\_adjudication\_log.csv} without editing the
original sheets. The pre-declared gate is Cohen's $\kappa \geq 0.70$ on the
verdict field~\cite{cohen1960}; if it is not met, the annotation manual is
rewritten from the disagreements and the pilot is re-run, which the protocol
states in advance is an expected outcome rather than a failure.

\paragraph{Current state.} Sheet \texttt{annotator\_B\_haaniah.csv} is complete:
all 15 cases carry a verdict, freshness state, review flag, reason code, the
evidence span identifiers, a written justification and a time estimate, and it
was committed in \texttt{8657955}. Sheet \texttt{annotator\_Abdullah.csv}
(\texttt{annotator\_A\_abdullah.csv}) is present but contains no verdicts yet.
Running the script today produces:

\begin{verbatim}
$ python3 docs/dataset/tools/agreement.py \
    docs/dataset/05_raw_annotations/annotator_A_abdullah.csv \
    docs/dataset/05_raw_annotations/annotator_B_haaniah.csv
No case labelled by both annotators yet.
$ echo $?
1
\end{verbatim}

\paragraph{Result.}
\textbf{[PENDING: agreement results --- run tools/agreement.py]}

The report will state, for the verdict, freshness and review-flag fields
separately: raw agreement, Cohen's kappa, the mean span Jaccard (IoU) and exact
span-match count, together with the list of disagreements handed to the
adjudicator. No kappa value is reported here because none has been computed;
estimating one would be the kind of claim this report is written to avoid.

% ---------------------------------------------------------------------------
\section{Acceptance Testing / Validation}
\label{sec:acceptance-testing}

Chapter~4 gives an acceptance check per functional requirement. Applied to the
current state of the repository, they sort into three groups:

\begin{longtable}{p{0.34\textwidth} p{0.16\textwidth} p{0.44\textwidth}}
\caption{Acceptance criteria groups and the validation status of each at the mid evaluation.}\label{tab:acceptance-groups}\\
\hline\textbf{Group} & \textbf{Requirements} & \textbf{Position} \\
\hline\endfirsthead
\hline\textbf{Group} & \textbf{Requirements} & \textbf{Position} \\
\hline\endhead
Partially demonstrated against mock data
  & FR-17, FR-11 (presentation half)
  & The change view loads and shows a known linked chain, criterion results,
    citations, verification passes and an audit trail --- from committed mock
    modules, not from a source system. Not acceptable as a final pass. \\
Cannot be tested yet (no implementation)
  & FR-01 to FR-10, FR-12 to FR-16, FR-18 to FR-29
  & Adapters, checkpoints, the graph, the analyser, the packet builder, the
    verifier, the gate, rules, routing and configuration do not exist. \\
Data-side validation complete
  & Dataset protocol (panel requirement)
  & Integrity checks pass on all 15 cases; annotation and adjudication are in
    progress; agreement pending. \\
\hline\end{longtable}

No requirement from Chapter~4 is declared \emph{met} in this report. The
honest status of the front-end rows is ``partial, against substituted data''.

% ---------------------------------------------------------------------------
\section{Edge and Failure Cases}
\label{sec:edge-cases}

\paragraph{Observed failure: the orphaned proof-of-concept module.}
\emph{Input.} \texttt{python3 -m pytest tests/} from \texttt{unify\_poc/}.
\emph{Why it matters.} The evaluation walkthrough will open the repository and
run something. A test directory that cannot collect is worse than no test
directory, because it implies coverage that does not exist, and it is exactly
the kind of gap the worksheet asks to be discussed rather than hidden.
\emph{Observed behaviour.} Collection aborts with
\texttt{ModuleNotFoundError: No module named 'unify'}; inspecting further, the
package directory in the tree holds a single module whose own imports
(\texttt{from .adapter import Adapter}, \texttt{from .normalizer import
FIGMA\_RE}) reference files deleted by the \texttt{066ef56} ``start fresh''
commit. \emph{Interpretation.} This establishes only that the surviving files
are not runnable; it says nothing about the quality of the code they contain,
and it does not affect the front end. \emph{Corrective action.} D-01: either
restore the two modules from history and make the test pass, or remove the
orphan before the evaluation tag. Owner: Abdullah. This is a decision to be
taken, not a fix that has been applied.

\paragraph{Designed failure case for the demo (not yet executable).}
The scenario matrix in Chapter~3 names T10 as the demonstration's failure case:
a fabricated citation injected into a verdict must be downgraded by the gate
and must never reach the reviewer as a decisive status. The gate that would
perform this check is designed (Chapter~5, Figure~\ref{fig:verifier-loop}) and
not implemented, so T10 cannot be run today. It is recorded here as planned, not
as passed. A substitute demonstration must not be presented in its place.

\paragraph{Edge case exercised by the build.}
The dynamic conflict route is the front end's only per-record path: if the mock
findings array were empty, \texttt{generateStaticParams()} would emit no detail
pages and links into the inbox would dead-end. The current build emits eight
pages, so the edge is visible in every build log as a page count rather than as
a silent omission.

% ---------------------------------------------------------------------------
\section{Defects and Corrective Actions}
\label{sec:defects}

\begin{longtable}{p{0.07\textwidth} p{0.30\textwidth} p{0.14\textwidth} p{0.15\textwidth} p{0.26\textwidth}}
\caption{Defects recorded during iteration 1 and the corrective action taken.}\label{tab:defects}\\
\hline\textbf{ID} & \textbf{Defect} & \textbf{Severity} & \textbf{Owner} &
  \textbf{Corrective action} \\
\hline\endfirsthead
\hline\textbf{ID} & \textbf{Defect} & \textbf{Severity} & \textbf{Owner} &
  \textbf{Corrective action} \\
\hline\endhead
D-01 & \texttt{unify\_poc} test module cannot be collected: package layout and
  two imported modules are missing from the tree. & High (credibility of the
  repository walkthrough) & Abdullah & Restore the modules from git history and
  make the suite pass, or delete the orphan before the evaluation tag. \\
D-02 & No CI: typecheck, lint, build and dataset checks run only when someone
  runs them. & Medium & Abdullah & Add a workflow running the three front-end
  commands and the integrity script on every pull request; require one other
  member's review. \\
D-03 & Agreement measurement blocked: annotator A's sheet has no verdicts. &
  Medium (blocks a panel requirement) & Abdullah & Complete the blind
  annotation, run \texttt{agreement.py}, then adjudicate; if
  $\kappa < 0.70$, revise the manual and re-pilot. \\
D-04 & No automated tests in \texttt{apps/web}; correctness rests on types, lint
  and the build. & Medium & Abdullah & Add component tests for the criterion
  result and resolution dock first, then the link-eligibility rules. \\
D-05 & Advisory: Node emits a \texttt{url.parse()} deprecation warning during
  build and lint, originating in framework dependencies. & Low & Abdullah &
  Track upstream; re-check on the next Next.js patch. \\
D-06 & The twelve diagrams are not stamped with the commit they were generated
  from, and are not regenerated at an evaluation tag. & Low & All three &
  Add the stamp to each source and regenerate at the tag (Chapter~3). \\
\hline\end{longtable}

No defect in this list has been closed. Reporting six open defects with owners
is preferable to reporting none, and every one of them is reproducible from the
commands above.

% ---------------------------------------------------------------------------
\section{Current Validation Status}
\label{sec:validation-status}

\begin{longtable}{p{0.20\textwidth} p{0.30\textwidth} p{0.24\textwidth} p{0.26\textwidth}}
\caption{Current validation status by requirement group.}\label{tab:validation-status}\\
\hline\textbf{Area} & \textbf{Current evidence} & \textbf{Limitation} &
  \textbf{Next validation action} \\
\hline\endfirsthead
\hline\textbf{Area} & \textbf{Current evidence} & \textbf{Limitation} &
  \textbf{Next validation action} \\
\hline\endhead
Functional correctness & Typecheck, lint and production build pass on all 19
  routes; 8 detail pages pre-rendered.
  & Only static checks exist; no behavioural tests, no back end.
  & Component tests (D-04), then scenario matrix T1--T11 once ingestion exists \\
Performance / resource use & Not applicable yet: no server, no model, nothing
  to time.
  & The hardware risk is completely unmeasured.
  & Run the local-inference spike (Chapter~3) and record latency, memory and
    quantisation \\
Model / algorithm quality & Not applicable yet: no model has been run, and no
  verifier exists.
  & Every verification claim in Chapters~4 and~5 is currently an unimplemented
    requirement.
  & Selection protocol first, then per-class verdict metrics on the frozen gold
    set \\
Evaluation data & 138 integrity checks pass on 15 cases from 3 repositories;
  annotation manual and dataset card written.
  & One annotator complete, one not started; no kappa; no frozen split; 15 cases
    is a pilot, not an oracle.
  & Finish annotation, compute agreement, adjudicate, then expand toward the
    development and gold sets \\
User / field validation & The interface was checked against its exported
  design source; no user has been observed using it.
  & No task-based session has taken place.
  & Run the crossover study against the manual baseline once verdicts come from
    the real pipeline \\
Security / privacy & No external requests at run time: fonts and icons are
  local, avatars are initials, no content leaves the browser.
  & The claims that matter (read-only scopes, encrypted tokens, egress
    inspection) belong to components that do not exist.
  & Egress-inspection test and credential handling review when the adapters are
    written \\
\hline\end{longtable}

The summary position is this: the artefacts that exist are checked and they
pass; the artefacts that carry the project's central claims do not exist yet, so
those claims are currently design commitments with acceptance criteria attached,
not measured results. Chapter~10 sequences the work that turns the right-hand
columns of the table above into the left-hand ones.
